\begingroup
\shorthandoff{"<>"}
\setlength{\arrayrulewidth}{1pt}
\arrayrulecolor{vlepsbased}
\begin{longtable}{|p{15cm}|}
  \hline
  \rowcolor{maincolor} 
  \multicolumn{1}{|l|}{\textcolor{white}{\textbf{Caso de Uso}}} \\
  \hline
  \endfirsthead
  \multicolumn{1}{c}{{\tablename\ \thetable{} -- Continuación}} \\
  \hline
  \rowcolor{maincolor} 
  \multicolumn{1}{|l|}{\textcolor{white}{\textbf{Caso de Uso (Continuación)}}} \\
  \hline
  \endhead
  \hline \multicolumn{1}{|r|}{\textit{Continúa en la siguiente página...}} \\ \hline
  \endfoot
  \hline
  \caption{Caso de Uso CU9 — Recuperación de contraseña olvidada.}
  \label{TB:CU9_RECUPERAR_PASSWORD}
  \endlastfoot

  \rowcolor{ulepsbased}
  \begin{tabularx}{\linewidth}{@{}X X@{}}
    \textbf{Identificador:} CU9 & \textbf{Actor principal:} Usuario no autenticado
  \end{tabularx} \\ \hline

  \textbf{Nombre:} Recuperación de contraseña olvidada \\ \hline

  \rowcolor{ulepsbased}
  \textbf{Descripción:} \newline 
  El usuario, al olvidar su contraseña, solicita un enlace temporal para poder crear una nueva clave y recuperar el acceso a su cuenta mediante el correo electrónico asociado. \\ \hline

  \textbf{Precondiciones:} \newline
  1. El sistema de envío de correos SMTP está configurado y operativo. \\ \hline

  \rowcolor{ulepsbased}
  \textbf{Flujo principal:} \newline
  1. El usuario accede a la opción "¿Olvidaste tu contraseña?" en la pantalla de acceso. \newline
  2. Introduce su dirección de correo electrónico y solicita la recuperación. \newline
  3. El sistema valida el formato e intenta localizar al usuario en la base de datos. \newline
  4. El sistema, por seguridad, devuelve una respuesta de confirmación idéntica independientemente de si el correo existe o no, para prevenir enumeración de cuentas. \newline
  5. En segundo plano, si el usuario existe, el sistema genera un token de un solo uso enlazado a dicho usuario y envía un correo electrónico. \newline
  6. El usuario accede a su bandeja de entrada, hace clic en el enlace, el cual le dirige a una vista especial en el frontend incluyendo los parámetros (UID y Token). \newline
  7. El usuario introduce una nueva contraseña segura y confirma. \newline
  8. El frontend envía la nueva clave al backend junto al UID y Token. \newline
  9. El sistema valida el token (que no haya expirado ni sido usado). \newline
  10. El sistema actualiza el hash de la contraseña en la base de datos. \newline
  11. Se informa al usuario y se le redirige al inicio de sesión. \\ \hline

  \textbf{Flujos alternativos:} \newline
  \textbf{FA-01:} Token inválido o expirado: \newline
  1. En el paso 9, el sistema detecta que el token ha expirado o que ya fue utilizado previamente. \newline
  2. Se rechaza la petición con un error HTTP 400. \newline
  3. El frontend muestra un mensaje indicando que el enlace ya no es válido y requiere solicitar uno nuevo. \\ \hline

  \rowcolor{ulepsbased}
  \textbf{Postcondiciones:} \newline
  1. La contraseña anterior es revocada y la nueva entra en funcionamiento. \newline
  2. El token de recuperación queda inmediatamente invalidado. \\ \hline

\end{longtable}
\endgroup
